%----------
%   IMPORTANTE
%----------

% Esta plantilla está basada en las recomendaciones de la guía "Trabajo fin de Máster: Escribir el TFM", que encontrarás en http://uc3m.libguides.com/TFM/escribir
% contiene recomendaciones de la Biblioteca basadas principalmente en estilos APA e IEEE, pero debes seguir siempre las orientaciones de tu Tutor de TFM y la normativa de TFM para tu titulación.



% ESTA PLANTILLA ESTÁ BASADA EN EL ESTILO IEEE


%----------
%	CONFIGURACIÓN DEL DOCUMENTO
%----------

\documentclass[12pt]{report} % fuente a 12pt

% MÁRGENES: 2,5 cm sup. e inf.; 3 cm izdo. y dcho.
\usepackage[
a4paper,
vmargin=2.5cm,
hmargin=3cm
]{geometry}

% INTERLINEADO: Estrecho (6 ptos./interlineado 1,15) o Moderado (6 ptos./interlineado 1,5)
\renewcommand{\baselinestretch}{1.15}
\parskip=6pt

% DEFINICIÓN DE COLORES para portada y listados de código
\usepackage[table]{xcolor}
\definecolor{azulUC3M}{RGB}{0,0,102}
\definecolor{gray97}{gray}{.97}
\definecolor{gray75}{gray}{.75}
\definecolor{gray45}{gray}{.45}

% Soporte para GENERAR PDF/A --es importante de cara a su inclusión en e-Archivo porque es el formato óptimo de preservación y a la generación de metadatos, tal y como se describe en http://uc3m.libguides.com/ld.php?content_id=31389625. 

% En la plantilla incluimos el archivo OUTPUT.XMPDATA. Puedes descargar este archivo e incluir los metadatos que se incorporarán al archivo PDF cuando compiles el archivo memoria.tex. Después vuelve a subirlo a tu proyecto. 
\usepackage[a-1b]{pdfx}

% ENLACES
\usepackage{hyperref}
\hypersetup{colorlinks=true,
	linkcolor=black, % enlaces a partes del documento (p.e. índice) en color negro
	urlcolor=blue, % enlaces a recursos fuera del documento en azul
	citecolor=black} 

% EXPRESIONES MATEMÁTICAS
\usepackage{amsmath,amssymb,amsfonts,amsthm}

% Codificación caracteres
\usepackage{txfonts} 
\usepackage[T1]{fontenc}
\usepackage[utf8]{inputenc}

% Definición idioma español
\usepackage[spanish, es-tabla]{babel} 
\usepackage[babel, spanish=spanish]{csquotes}
\AtBeginEnvironment{quote}{\small}

% diseño de PIE DE PÁGINA
\usepackage{fancyhdr}
\pagestyle{fancy}
\fancyhf{}
\renewcommand{\headrulewidth}{0pt}
\rfoot{\thepage}
\fancypagestyle{plain}{\pagestyle{fancy}}

% DISEÑO DE LOS TÍTULOS de las partes del trabajo (capítulos y epígrafes o subcapítulos)
\usepackage{titlesec}
\usepackage{titletoc}
\titleformat{\chapter}[block]
{\large\bfseries\filcenter}
{\thechapter.}
{5pt}
{\MakeUppercase}
{}
\titlespacing{\chapter}{0pt}{0pt}{*3}
\titlecontents{chapter}
[0pt]                                               
{}
{\contentsmargin{0pt}\thecontentslabel.\enspace\uppercase}
{\contentsmargin{0pt}\uppercase}                        
{\titlerule*[.7pc]{.}\contentspage}                 

\titleformat{\section}
{\bfseries}
{\thesection.}
{5pt}
{}
\titlecontents{section}
[5pt]                                               
{}
{\contentsmargin{0pt}\thecontentslabel.\enspace}
{\contentsmargin{0pt}}
{\titlerule*[.7pc]{.}\contentspage}

\titleformat{\subsection}
{\normalsize\bfseries}
{\thesubsection.}
{5pt}
{}
\titlecontents{subsection}
[10pt]                                               
{}
{\contentsmargin{0pt}                          
	\thecontentslabel.\enspace}
{\contentsmargin{0pt}}                        
{\titlerule*[.7pc]{.}\contentspage}  


% DISEÑO DE TABLAS
\usepackage{multirow} % permite combinar celdas 
\usepackage{caption} % para personalizar el título de tablas y figuras
\usepackage{floatrow} % utilizamos este paquete y sus macros \ttabbox y \ffigbox para alinear los nombres de tablas y figuras de acuerdo con el estilo definido.
\usepackage{array} % con este paquete podemos definir en la siguiente línea un nuevo tipo de columna para tablas: ancho personalizado y contenido centrado
\newcolumntype{P}[1]{>{\centering\arraybackslash}p{#1}}
\DeclareCaptionFormat{upper}{#1#2\uppercase{#3}\par}

% Diseño de tabla para ingeniería
\captionsetup*[table]{
	format=upper,
	name=TABLA,
	justification=centering,
	labelsep=period,
	width=.75\linewidth,
	labelfont=small,
	font=small
}

% DISEÑO DE FIGURAS. 
\usepackage{graphicx}
\graphicspath{{imagenes/}} % ruta a la carpeta de imágenes

% Diseño de figuras para ingeniería
\captionsetup[figure]{
	format=hang,
	name=Fig.,
	singlelinecheck=off,
	labelsep=period,
	labelfont=small,
	font=small		
}

% NOTAS A PIE DE PÁGINA
\usepackage{chngcntr} % para numeración continua de las notas al pie
\counterwithout{footnote}{chapter}

% LISTADOS DE CÓDIGO
% soporte y estilo para listados de código. Más información en https://es.wikibooks.org/wiki/Manual_de_LaTeX/Listados_de_código/Listados_con_listings
\usepackage{listings}

% definimos un estilo de listings
\lstdefinestyle{estilo}{ frame=Ltb,
	framerule=0pt,
	aboveskip=0.5cm,
	framextopmargin=3pt,
	framexbottommargin=3pt,
	framexleftmargin=0.4cm,
	framesep=0pt,
	rulesep=.4pt,
	backgroundcolor=\color{gray97},
	rulesepcolor=\color{black},
	%
	basicstyle=\ttfamily\footnotesize,
	keywordstyle=\bfseries,
	stringstyle=\ttfamily,
	showstringspaces = false,
	commentstyle=\color{gray45},     
	%
	numbers=left,
	numbersep=15pt,
	numberstyle=\tiny,
	numberfirstline = false,
	breaklines=true,
	xleftmargin=\parindent
}

\captionsetup*[lstlisting]{font=small, labelsep=period}
% fijamos el estilo a utilizar 
\lstset{style=estilo}
\renewcommand{\lstlistingname}{\uppercase{Código}}


%BIBLIOGRAFÍA 

% CONFIGURACIÓN PARA LA BIBLIOGRAFÍA IEEE
\usepackage[backend=biber, style=ieee, isbn=false,sortcites, maxbibnames=6, minbibnames=1]{biblatex} % Configuración para el estilo de citas de IEEE, recomendado para el área de ingeniería. "maxbibnames" indica que a partir de 6 autores trunque la lista en el primero (minbibnames) y añada "et al." tal y como se utiliza en el estilo IEEE.

% Añadimos las siguientes indicaciones para mejorar la adaptación del estilos en español
\DefineBibliographyStrings{spanish}{%
	andothers = {et\addabbrvspace al\adddot}
}
\DefineBibliographyStrings{spanish}{
	url = {\adddot\space[En línea]\adddot\space Disponible en:}
}
\DefineBibliographyStrings{spanish}{
	urlseen = {Acceso:}
}
\DefineBibliographyStrings{spanish}{
	pages = {pp\adddot},
	page = {p.\adddot}
}

\addbibresource{referencias.bib} % llama al archivo referencias.bib en el que deberá estar la bibliografía utilizada


%-------------
%	DOCUMENTO
%-------------

\begin{document}
\pagenumbering{roman} % Se utilizan cifras romanas en la numeración de las páginas previas al cuerpo del trabajo
	
%----------
%	PORTADA
%----------	
\begin{titlepage}
	\begin{sffamily}
	\color{azulUC3M}
	\begin{center}
		\begin{figure}[H] %incluimos el logotipo de la Universidad
			\makebox[\textwidth][c]{\includegraphics[width=16cm]{logo_UC3M.png}}
		\end{figure}
		\vspace{2.5cm}
		\begin{Large}
			Máster Universitario en Ingeniería Industrial\\			
			 2023-2024\\ %Indica el curso académico
			\vspace{2cm}		
			\textsl{Trabajo Fin de Máster}
			\bigskip
			
		\end{Large}
		 	{\Huge ``Título del trabajo''}\\
		 	\vspace*{0.5cm}
	 		\rule{10.5cm}{0.1mm}\\
			\vspace*{0.9cm}
			{\LARGE Antonio Pérez García}\\ 
			\vspace*{1cm}
		\begin{Large}
			Tutor\\
			Gabriel Lázaro\\
			Lugar y fecha de presentación prevista\\
		\end{Large}
	\end{center}
	\vfill
	\color{black}
	\fbox{
		\begin{minipage}{\linewidth}
		\textbf{DETECCIÓN DEL PLAGIO}\\
		\footnotesize{La Universidad utiliza el programa \textbf{Turnitin Feedback Studio} para comparar la originalidad del trabajo entregado por cada estudiante con millones de recursos electrónicos y detecta aquellas partes del texto copiadas y pegadas. Copiar o plagiar en un TFM es considerado una \textbf{\underline{Falta Grave}}, y puede conllevar la expulsión definitiva de la Universidad.}\end{minipage}}
	
	% SI NUESTRO TRABAJO SE VA A PUBLICAR CON UNA LICENCIA CREATIVE COMMONS, INCLUIR ESTAS LÍNEAS. ES LA OPCIÓN RECOMENDADA.
	\noindent\includegraphics[width=4.2cm]{creativecommons.png}\\ %incluimos el logotipo de Creative Commons
	\footnotesize{Esta obra se encuentra sujeta a la licencia Creative Commons \textbf{Reconocimiento - No Comercial - Sin Obra Derivada}}
	
	\end{sffamily}
\end{titlepage}

\newpage %página en blanco o de cortesía
\thispagestyle{empty}
\mbox{}

%----------
%	RESUMEN Y PALABRAS CLAVE
%----------	
\renewcommand\abstractname{\large\bfseries\filcenter\uppercase{Resumen}}
\begin{abstract}
\thispagestyle{plain}
\setcounter{page}{3}
	
	% ESCRIBIR EL RESUMEN AQUÍ
	
	\textbf{Palabras clave:}
	% Escribir las palabras clave aquí
	
	\vfill
\end{abstract}

	\newpage % página en blanco o de cortesía
	\thispagestyle{empty}
	\mbox{}


%----------
%	DEDICATORIA
%----------	
\chapter*{Dedicatoria} % \chapter* evita que aparezca en el índice

\setcounter{page}{5}
	
	% ESCRIBIR LA DEDICATORIA AQUÍ	
		
	\vfill
	
	\newpage % página en blanco o de cortesía
	\thispagestyle{empty}
	\mbox{}
	

%----------
%	ÍNDICES
%----------	

%--
% Índice general
%-
\tableofcontents
\thispagestyle{fancy}

\newpage % página en blanco o de cortesía
\thispagestyle{empty}
\mbox{}

%--
% Índice de figuras. Si no se incluyen, comenta las líneas siguientes
%-
\listoffigures
\thispagestyle{fancy}

\newpage % página en blanco o de cortesía
\thispagestyle{empty}
\mbox{}

%--
% Índice de tablas. Si no se incluyen, comenta las líneas siguientes
%-
\listoftables
\thispagestyle{fancy}

\newpage % página en blanco o de cortesía
\thispagestyle{empty}
\mbox{}


%----------
%	MEMORIA
%----------	
\clearpage
\pagenumbering{arabic} % numeración con números arábigos para el resto de la memoria.	

\chapter{Introducción}\label{cap1}

La previsión de demanda es una herramienta importante para la planificación estratégica de múltiples sectores como el comercio minorista, la producción industrial y la logística. En el sector del \textit{aftermarket} del automóvil, la capacidad de predecir con la suficiente precisión la demanada de productos es fundamental para optimizar la cadena de suministro, reducir los costes de inventario y mejorar el servicio al cliente. La variabilidad de la demanda debida a factores como la estacionalidad, las tendencias económicas, las normativas medioambientales y los cambios en el comportamiento de los consumidores hace que las empresas se enfrenten a desafíos considerables a la hora de planificar su cadena de suministros de forma eficiente.

Las predicciones se pueden realizar mediante métodos intuitivos, en base a la razón de los dirigientes, o mediante métodos cuantitativos, utilizando datos históricos. Tradicionalmente, se han utilizado métodos estadísticos convencionales para realizar predicciones, como pueden ser modelos de media móvil, suavizado exponencial o el método de Holt-Winters \cite{winters}. Sin embargo, los avances en el campo de la inteligencia artificial ha permitido el desarrollo de modelos de aprendizaje automático que son capaces de interpretar patrones complejos en el histórico de datos. Además, estos algoritmos pueden integrar variables externas que impacten en los datos de estudio, como puede ser el precio del combustible o la meteorología.

En este contexto, el presente trabajo se centra en la aplicación de modelos de aprendizaje automático utilizando los datos de ventas de varias familias de producto de una empresa de recambios de coche situada en Andalucía, con tiendas en las zonas de la sierra limítrofes entre Cádiz y Málaga. Para ello, se utiliza el lenguaje de programación Python \cite{python}, uno de los más utilizados en el campo de la ciencia de datos, junto con la librería Darts \cite{darts}, especializada en \textit{forecasting} en series temporales.

\section{Estructura del trabajo}

En el Capítulo \ref{cap2}: Marco teórico, se realiza una revisión de los algoritmos de predicción de series temporales más relevantes para el caso de estudio y una descripción general del ámbito del \textit{forecast}.

A continuación, en el Capítulo \ref{cap3}: Metodología, se detallan las series históricas de datos utilizados, los algoritmos implementados y así como las métricas para la evaluación de los resultados.

En el Capítulo \ref{cap4}: Desarrollo, se describe la implementación y el código utilizado para el procesamiento de los datos, el entrenamiento de los modelos y las predicciones realizadas.

Una vez realizadas las predicciones, se estudian los resultados obtenidos incluyendo las gráficas y métricas en el Capítulo \ref{cap5}: Resultados.

Para finalizar el trabajo, se incluye el Capítulo \ref{cap6}: Conclusiones, donde se realizan las reflexiones finales del estudio se sugieren propuestas para proyectos futuros.


% BLOQUE A1: NIVEL TECNOLÓGICO

% BLOQUE A2: PROBLEMA DE LA INDUSTRIA A RESOLVER

% Relevando de la acción de repuestos en actividad a soporte en servicios

% Caso de uso<<<<<<<<<<<<<<<<<<<ss

% BLOQUE B: DESARROLLO DE LA SOLUCIÓN

% BLOQUE C: COMO LA SOLUCIÓN HA SOLUCIONADO EL PROBLEMA Y LÍNEAS DE TRABAJO FUTURA

\chapter{Marco teórico}\label{cap2}

El \textit{forecasting} es la predicción uno o varios eventos futuros \cite{intro}. Esta simple descripción permite que un ámbito de investigación tenga impacto en diversos sectores como los mercados, las finanzas o la meteorología. Tanto la predicción del valor de una acción en los próximos minutos en bolsa como la predicción de las ventas de una tienda durante el próximo trimestre se benefician de los avances en este ámbito. Las predicciones pueden ser a corto plazo, medio y largo plazo, por lo que pueden ser utilizadas para optimizar las decisiones operativas, tácticas y estratégicas de las empresas. En el contexto empresarial, una buena previsión permite anticiparse a cambios en la demanda ajustando la producción, el suministro y la distribución para maximizar la rentabilidad y minimizar los riesgos asociados con el exceso o falta de stock. 

\section{Métodos tradicionales}

\section{Aprendizaje automático}


\section{Modelos}

Los modelos matemáticos son herramientas matemáticas o estadísticas diseñadas para identificar patrones en las series de datos con el fin de predecir valores futuros. Cada modelo tiene una arquitectura que captura diferentes características y distintos tipos de patrones de los datos, como la tendencia, la estacionalidad y las autocorrelaciones, utilizando dicha información para generar predicciones basadas en el comportamiento pasado de la serie temporal.

\subsection{Prophet}

Prophet \cite{prophet} es una herramienta desarrollada por Facebook (ahora conocida como Meta) para el análisis de series temporales. Está especializada en datos que presentan patrones estacionales y tendencia. Está orientado al usuario final sin requerir configuración compleja específica para su aplicación. Se utiliza como librería en Python o R \cite{R}. El modelo incluido descompone la serie temporal en tendencia y estacionalidad. Además, permite añadir los festivos locales para mejorar las predicciones.

Se utiliza un modelo en el que la tendencia se modela de forma flexible, permitiendo cambios no lineales mediante \textit{changepoints}. Evita los valores atípicos y los datos faltantes, dando robustez al modelo frente a series temporales ruidosas o incompletas.

Su uso en este trabajo es de \textit{plug and play}, ya que se instancia la librería por defecto, a excepción de la especificación de los festivos de España (la librería no incluye la posiblidad de añadir los festivos de Andalucía).

Es necesario destacar que es el único modelo de los implementados que no depende de PyTorch \cite{pytorch}.

\subsection{RNN}
\subsection{N-BEATS y N-HiTS}
\subsection{TCN}
\subsection{DLinear y NLinear}
\subsection{TiDE}
\subsection{TSMixer}

	% COMENZAR A ESCRIBIR la MEMORIA
	
	% IMPORTANTE: en LaTeX hay una serie de caracteres especiales, que son: # $ % & \ ^ _ { } ~. Si aparecen en el texto, tendrás un error al compilar. La mayoría se pueden escapar escribiendo \ delante. Para \ utiliza \textbackslash ; para ^ \textasciitilde y para ~ \textasciicircum.

    % COMO incluir una FIGURA siguiendo las recomendaciones de la Guía: Alineación del título: izquierda, en la parte inferior de la figura; fuente a 10pt (el resto de la memoria está a 12); Numeración de la figura: se identificarán con números arábigos consecutivos tras la palabra Tabla. Si se sigue un esquema numérico de capítulos, el número de la tabla debe identificar con su primer dígito al capítulo, seguido por un punto y el número consecutivo que corresponda (Tabla 1.1, etc.); Propiedad intelectual: Se debe indicar la fuente de origen de la información en la parte inferior de la figura, a continuación del título.
    
    % EJEMPLO DE INCLUSIÓN de una figura:
    % \begin{figure}[H]
    % 	\ffigbox[\FBwidth] {
    % 	\caption[Nombre que aparecerá en el índice]{Nombre que aparecerá debajo de la figura}
    % 	}
    % 	{\includegraphics[scale=0.6]{imagenes/creativecommons.png}}
    % \end{figure}
    
    % COMO incluir una TABLA siguiendo las recomendaciones de la Guía: Alineación del título: Centrado en la parte superior de de la tabla, en la parte superior de la tabla; fuente a 10pt (el resto de la memoria está a 12) y en mayúsculas; Numeración: se identificarán con números arábigos consecutivos tras la palabra Tabla. Si se sigue un esquema numérico de capítulos, el número de la tabla debe identificar con su primer dígito al capítulo, seguido por un punto y el número consecutivo que corresponda (Tabla 1.1, etc.); Propiedad intelectual: Se debe indicar la fuente de origen de la información en la parte inferior de la tabla.
    
    % EJEMPLO DE INCLUSIÓN de una tabla:
% \begin{table}[H]
% 	\ttabbox[\FBwidth]
% 	{\caption{Lorem ipsum}}
% 	{\begin{tabular}{|c|P{1.5cm}|c|P{1.5cm}|P{2cm}|c|P{1.5cm}|P{2cm}|}
% 		\hline
% 		\multicolumn{2}{|c|}{\textbf{I}} & \multicolumn{2}{c|}{\textbf{II}} & \multicolumn{3}{c|}{\textbf{III}} & \textbf{IV} \\
% 		\hline
% 		x & y & x & y & x & y & x & y \\
% 		\hline
% 		10.0 & 8.04 & 10.0 & 9.14 & 10.0 & 7.46 & 8.0 & 6.58 \\
% 		\hline
% 		8.0 & 6.95 & 8.0 & 8.14 & 8.0 & 6.77 & 8.0 & 5.76 \\
% 		\hline
% 		13.0 & 7.58 & 13.0 & 8.74 & 13.0 & 12.74 & 8.0 & 7.71 \\
% 		\hline
% 		9.0 & 8.81 & 9.0 & 8.77 & 9.0 & 7.11 & 8.0 & 8.84 \\
% 		\hline
% 		11.0 & 8.33 & 11.0 & 9.26 & 11.0 & 7.81 & 8.0 & 8.47 \\
% 		\hline
% 		14.0 & 9.96 & 14.0 & 8.10 & 14.0 & 8.84 & 8.0 & 7.04 \\
% 		\hline
% 		6.0 & 7.24 & 6.0 & 6.13 & 6.0 & 6.08 & 8.0 & 5.25 \\
% 		\hline
% 		4.0 & 4.26 & 4.0 & 3.10 & 4.0 & 5.39 & 19.0 & 12.50 \\
% 		\hline
% 		12.0 & 10.84 & 12.0 & 9.13 & 12.0 & 8.15 & 8.0 & 5.56 \\
% 		\hline
% 		7.0 & 4.82 & 7.0 & 7.26 & 7.0 & 6.42 & 8.0 & 7.91 \\
% 		\hline
% 		5.0 & 5.68 & 5.0 & 4.74 & 5.0 & 5.73 & 8.0 & 6.89 \\
% 		\hline
% 		\multicolumn{5}{l}{Fuente: BOE}
% 	\end{tabular}}
% \end{table}


%----------
%	BIBLIOGRAFÍA
%----------	

%\nocite{*} % Si quieres que aparezcan en la bibliografía todos los documentos que la componen (también los que no estén citados en el texto) descomenta está línea

\clearpage
\addcontentsline{toc}{chapter}{Bibliografía}
\setquotestyle[british]{english} % Cambiamos el tipo de cita porque en el estilo IEEE se usan las comillas inglesas.
\printbibliography



%----------
%	ANEXOS
%----------	

% Si tu trabajo incluye anexos, puedes descomentar las siguientes líneas
%\chapter* {Anexo x}
%\pagenumbering{gobble} % Las páginas de los anexos no se numeran



\end{document}